\section{Restricting the exponents}
\label{sec:exponent_restriction}

We now pass from real exponents to an arbitrary additive submonoid
$\Gamma\subseteq\mathbb{R}$.  Although extension of exponents gives an
injective map
\[
  A((t^\Gamma))\longrightarrow A((t^{\mathbb{R}})),
\]
we do not assume that this map is faithfully flat.  Instead, we first
construct an injective map from a descent datum to a constant object and then
apply the same construction to its dual.  Finite-projective duality supplies a
right inverse.

\begin{definition}[Extension and restriction of exponents]
  \label{def:exponent_extension}
  \lean{Novikov.HasExponentSupport,
    Novikov.extendExponents,
    Novikov.restrictExponents,
    Novikov.mem_range_extendExponents_iff,
    Novikov.extendExponentsRingHom}
  \leanok
  \uses{def:novikov_series, thm:novikov_ring}
  For variables $t_1,\ldots,t_n$, extension by zero is the map
  \[
    \iota_*\colon A((t_1^\Gamma,\dotsc,t_n^\Gamma))
      \longrightarrow A((t_1^{\mathbb{R}},\dotsc,t_n^{\mathbb{R}})).
  \]
  Its coefficient at $d=(d_1,\ldots,d_n)\in\mathbb{R}^n$ is the coefficient
  of the original series when every $d_i$ lies in $\Gamma$, and is zero
  otherwise.  It is an injective ring homomorphism.

  A real-exponent series has \emph{$\Gamma$-support} if every coordinate of
  every exponent in its support lies in $\Gamma$.  The image of $\iota_*$ is
  exactly the set of series with $\Gamma$-support.  Restriction of coefficients
  to $\Gamma^n$ is inverse to extension on this set.
\end{definition}

\begin{proposition}[Exponent base change]
  \label{prop:exponent_base_change}
  \lean{Novikov.Descent.extendExponents_substitute_of_injective,
    Novikov.Descent.exponentInclusionCHom,
    Novikov.Descent.exponentInclusionExtendedCHom,
    Novikov.Descent.exponentBaseChangeFunctor,
    Novikov.Descent.exponentConstantBaseChangeIso}
  \leanok
  \uses{def:exponent_extension, def:novikov_descent_datum,
    def:constant_descent_datum}
  Extension by zero commutes with substitution along each face map in the
  Novikov cosimplicial ring.  It therefore defines a base-change functor
  \[
    (-)_{\mathbb{R}}\colon
      \mathsf{NovDD}(A,\Gamma)\longrightarrow
      \mathsf{NovDD}(A,\mathbb{R}).
  \]
  For every finite projective $A$-module $P$, there is a canonical
  isomorphism
  \[
    \bigl(\operatorname{Const}_\Gamma(P)\bigr)_{\mathbb{R}}
      \cong \operatorname{Const}_{\mathbb{R}}(P).
  \]
\end{proposition}

\begin{proof}
  \leanok
  Every face map is substitution along an injective map of finite variable
  sets.  A coefficient calculation shows that extending before or after such
  a substitution gives the same series.  The level-one, level-two, and
  level-three extension homomorphisms thus form a morphism of cosimplicial
  rings.  At level zero the map is the identity on $A$, so abstract constant
  base change, followed by cancellation of $A\otimes_A P$, gives the displayed
  comparison for constant objects.
\end{proof}

\begin{lemma}[Support of a real trivialization]
  \label{lem:real_trivialization_exponent_support}
  \lean{Novikov.Descent.realTrivializedSeries,
    Novikov.Descent.realTrivializedSeries_has_exponent_support}
  \leanok
  \uses{prop:exponent_base_change, def:constant_descent_datum}
  Let $M\in\mathsf{NovDD}(A,\Gamma)$, let $P$ be finite projective over $A$,
  and suppose that
  \[
    e\colon M_{\mathbb{R}}\xrightarrow{\sim}
      \operatorname{Const}_{\mathbb{R}}(P)
  \]
  is an isomorphism.  Identify the underlying module of the constant object
  with the module of $P$-valued real Novikov series.  For $m\in M$, let
  $y_m$ be the series corresponding to $e(1\otimes m)$.  Then
  \[
    \operatorname{supp}(y_m)\subseteq\Gamma.
  \]
\end{lemma}

\begin{proof}
  \leanok
  Realize the two pullbacks of $M$ as two-variable, $P$-valued real Novikov
  series.  On a pure tensor, the first realization has the form
  \[
    a\otimes m\longmapsto
      \iota_*(a)\,\operatorname{sub}_1(y_m).
  \]
  Here $\iota_*(a)$ has support in $\Gamma\times\Gamma$, while
  $\operatorname{sub}_1(y_m)$ has support in $\mathbb{R}\times\{0\}$.
  Consequently every element in the image of the first realization is
  supported in $\mathbb{R}\times\Gamma$.

  The descent isomorphism identifies the second realization of $1\otimes m$
  with an element in this first image.  The second realization substitutes
  $y_m$ into the second coordinate.  Evaluating its coefficient at $(0,d)$
  therefore shows that every nonzero coefficient of $y_m$ has
  $d\in\Gamma$.
\end{proof}

\begin{proposition}[Restricting a real trivialization]
  \label{prop:restricted_real_trivialization}
  \lean{Novikov.Descent.restrictedTrivializedSeries,
    Novikov.Descent.restrictRealTrivialization,
    Novikov.Descent.restrictRealTrivialization_injective,
    Novikov.Descent.restrictRealTrivialization_baseChange}
  \leanok
  \uses{lem:real_trivialization_exponent_support,
    prop:exponent_base_change}
  Under the hypotheses of
  Lemma~\ref{lem:real_trivialization_exponent_support}, coefficient restriction
  defines a morphism of descent data
  \[
    g\colon M\longrightarrow\operatorname{Const}_\Gamma(P).
  \]
  This morphism is injective.  After extension to real exponents and the
  canonical comparison for constant objects, it recovers $e$.
\end{proposition}

\begin{proof}
  \leanok
  Define $g(m)$ by restricting the coefficients of $y_m$ to $\Gamma$ and then
  using the standard identification of $P$-valued Novikov series with the
  underlying module of $\operatorname{Const}_\Gamma(P)$.  The preceding lemma
  says that extending this restricted series recovers $y_m$.

  Restriction of an arbitrary real series need not be
  $A((t^\Gamma))$-linear when $\Gamma$ is only a submonoid.  Linearity of $g$
  is instead proved after applying the injective extension map, where it
  reduces to linearity of $e$.  The same strategy proves compatibility with
  the descent isomorphisms: after real exponent base change the desired square
  is the square for $e$, and equality is reflected back through the
  level-two extension map using finite projectivity.  Finally, if
  $g(m)=g(n)$, extension gives $e(1\otimes m)=e(1\otimes n)$.  Injectivity of
  $m\mapsto1\otimes m$ for a finite projective module over an injective ring
  map gives $m=n$.
\end{proof}

\begin{proposition}[Duality makes the restricted map invertible]
  \label{prop:restricted_real_trivialization_iso}
  \lean{Novikov.Descent.novikovConstantDualIso,
    Novikov.Descent.realDualTrivialization,
    Novikov.Descent.restrictRealDualTrivialization,
    Novikov.Descent.fromRestrictedDual,
    Novikov.Descent.restrictRealTrivialization_comp_fromRestrictedDual,
    Novikov.Descent.restrictRealTrivialization_bijective,
    Novikov.Descent.restrictRealTrivializationIso}
  \leanok
  \uses{prop:restricted_real_trivialization,
    def:constant_descent_datum}
  The morphism $g$ of
  Proposition~\ref{prop:restricted_real_trivialization} is an isomorphism:
  \[
    M\xrightarrow{\sim}\operatorname{Const}_\Gamma(P).
  \]
\end{proposition}

\begin{proof}
  \leanok
  Finite projectivity identifies a module with its double dual and gives a
  canonical comparison
  \[
    c\colon\operatorname{Const}_\Gamma(P^\vee)
      \xrightarrow{\sim}\operatorname{Const}_\Gamma(P)^\vee.
  \]
  Base change of duals is canonically the dual of base change.  Applying these
  comparisons to the dual of $e$, and then applying
  Proposition~\ref{prop:restricted_real_trivialization}, produces an injective
  morphism
  \[
    h\colon M^\vee\longrightarrow
      \operatorname{Const}_\Gamma(P^\vee).
  \]

  The primal and dual constructions are compatible.  After extending
  exponents this follows from the naturality of the base-change/dual
  comparison and the fact that $g_{\mathbb{R}}$ is $e$.  Reflecting the equality
  back to $\Gamma$ gives
  \[
    h\circ g^\vee=c^{-1}.
  \]
  Dualize $h$, compose with $c$ and the double-dual evaluations, and obtain
  \[
    q\colon\operatorname{Const}_\Gamma(P)\longrightarrow M.
  \]
  Pairing $g(q(p))$ with an arbitrary element of
  $\operatorname{Const}_\Gamma(P)^\vee$ and using the displayed identity shows
  that
  \[
    g\circ q=\operatorname{id}.
  \]
  Thus $g$ is surjective; it was already injective by
  Proposition~\ref{prop:restricted_real_trivialization}.  Its underlying
  linear map is therefore a linear equivalence, and the descent compatibility
  packages it as the claimed isomorphism.
\end{proof}

\begin{theorem}[Theorem~4.9: the main Novikov descent theorem]
  \label{thm:novikov_descent_equivalence}
  \lean{Novikov.Descent.vectToNovikovDescent_essSurj,
    Novikov.Descent.vectToNovikovDescent_isEquivalence,
    Novikov.Descent.vectToNovikovDescent_equivalence}
  \leanok
  \uses{thm:novikov_descent_equivalence_real,
    prop:restricted_real_trivialization_iso,
    prop:vect_to_novikov_descent_fully_faithful}
  Let $A$ be a commutative ring and let
  $\Gamma\subseteq\mathbb{R}$ be an additive submonoid.  The constant-descent
  functor
  \[
    \mathsf{Vect}(A)\longrightarrow\mathsf{NovDD}(A,\Gamma)
  \]
  is an equivalence of categories.
\end{theorem}

\begin{proof}
  \leanok
  Let $M\in\mathsf{NovDD}(A,\Gamma)$.  Extend its exponents to obtain
  $M_{\mathbb{R}}$.  The real-exponent theorem gives a finite projective
  $A$-module $P$ and an isomorphism
  \[
    M_{\mathbb{R}}\cong\operatorname{Const}_{\mathbb{R}}(P).
  \]
  Proposition~\ref{prop:restricted_real_trivialization_iso} restricts this to
  an isomorphism
  $M\cong\operatorname{Const}_\Gamma(P)$, proving essential surjectivity.
  Full faithfulness is
  Proposition~\ref{prop:vect_to_novikov_descent_fully_faithful}.
\end{proof}
